% ============================================================
% Report: Flow-Matching Speech Enhancement via
%         Dynamic Multi-Layer Tokenizer Alignment
% CVPR format, ≤5 pages, two columns
% ============================================================
\documentclass[10pt,twocolumn,letterpaper]{article}
\usepackage[final]{cvpr}
\usepackage{times}
\usepackage[utf8]{inputenc}
\usepackage[T1]{fontenc}
\usepackage{amsmath,amssymb}
\usepackage{graphicx}
\usepackage{booktabs}
\usepackage{hyperref}
\usepackage[capitalize]{cleveref}
\usepackage{enumitem}
\usepackage{caption}
\usepackage{subcaption}
\usepackage{microtype}
\usepackage{xcolor}
\usepackage{float}

\hypersetup{colorlinks=true,linkcolor=blue!60!black,citecolor=blue!60!black,urlcolor=blue!60!black}

\graphicspath{{../poster/images/}}

\title{Flow-Matching Speech Enhancement via Dynamic Multi-Layer Tokenizer Alignment}

\author{
  Yiheng Chen \quad Andr\'{e} Dal Bosco\\
  \'{E}cole Polytechnique\\
  CSC\_52002\_EP --- IA G\'{e}n\'{e}rative Multimodale (2025--2026)\\
  Teachers: Vicky Kalogeiton, Xi Wang\\
  {\tt\small \href{https://github.com/VictorChen2002/Speech-Enhancement-Project}{github.com/VictorChen2002/Speech-Enhancement-Project}}
}

\begin{document}
\maketitle

% ============================================================
\begin{abstract}
We present a flow-matching speech enhancement system that operates in the continuous latent space of two pre-trained audio tokenizers.
A Rectified Flow model learns a direct velocity field from noisy DAC latents to clean DAC latents, conditioned on features extracted from the MOSS-Audio-Tokenizer.
Our core contribution is a \emph{multi-layer conditioning} strategy: instead of using only the last hidden layer of the MOSS tokenizer (standard practice), we extract all 32 hidden-layer outputs from the deepest encoder stage and fuse them via a learnable weighted sum.
We further propose a \emph{time-dependent fusion} variant where a lightweight MLP maps the ODE timestep to per-layer weights, allowing different denoising stages to emphasise different feature types.
A controlled ablation across four conditioning strategies shows consistent, monotonic improvements in PESQ, STOI, and FAD, confirming that multi-layer tokenizer features provide complementary information for speech enhancement.
\end{abstract}

% ============================================================
\section{Introduction}
\label{sec:intro}

Speech enhancement---the task of extracting clean speech from a noisy recording---is fundamental to applications including automatic speech recognition, hearing aids, and voice communication.
Classical approaches operate in the time-frequency domain with supervised regression~\cite{lu2013speech}, while recent work has explored generative models that produce natural-sounding speech by modelling the distribution of clean signals~\cite{richter2023speech}.

A promising direction is latent-space generative modelling: rather than operating on raw waveforms or spectrograms, one encodes audio into a compact continuous latent space using a pre-trained codec and performs generation there.
This dramatically reduces the sequence length and allows the generative model to focus on high-level structure.

In this project, we adopt \textbf{Rectified Flow}~\cite{lipman2023flow}---a flow-matching framework that learns a straight-line velocity field between two distributions---to map noisy audio latents to clean audio latents.
A \textbf{Diffusion Transformer} (DiT)~\cite{peebles2023scalable} serves as the velocity predictor, conditioned via cross-attention on features of the \textbf{MOSS-Audio-Tokenizer}~\cite{gong2026moss}, a large-scale (1.6B parameter) hierarchical audio encoder.

\paragraph{Core contribution.}
Pre-trained feature extractors contain rich information distributed across their layers: early layers capture fine-grained acoustic details while deeper layers encode more abstract semantic and prosodic features.
In standard practice, only the \emph{last} hidden layer is used for conditioning.
We hypothesise that leveraging \emph{all} layers can improve enhancement quality.
We design two fusion strategies---static and time-dependent---and conduct a systematic four-way ablation to isolate the effect of multi-layer conditioning.

% ============================================================
\section{Background}
\label{sec:background}

\subsection{Rectified Flow}

Rectified Flow~\cite{lipman2023flow,liu2023flow} defines a generative process by learning a velocity field that transports samples along straight paths between a source and target distribution.
Given paired samples $\mathbf{x}_0 \sim p_{\text{source}}$ and $\mathbf{x}_1 \sim p_{\text{target}}$, the interpolation at time $t \in [0,1]$ is:
\begin{equation}
  \mathbf{x}_t = t\,\mathbf{x}_1 + (1 - t)\,\mathbf{x}_0
  \label{eq:interpolation}
\end{equation}
The target velocity field is constant along each path: $\mathbf{v}^* = \mathbf{x}_1 - \mathbf{x}_0$.
A neural network $\mathbf{v}_\theta$ is trained via:
\begin{equation}
  \mathcal{L} = \mathbb{E}_{t \sim U(0,1)}\!\left[\left\|\mathbf{v}_\theta(\mathbf{x}_t, t, \mathbf{c}) - (\mathbf{x}_1 - \mathbf{x}_0)\right\|^2\right]
  \label{eq:loss}
\end{equation}
In our setting $\mathbf{x}_0$ is \emph{not} Gaussian noise but the noisy DAC latent---the flow learns a denoising transformation starting from meaningful corrupted input.

\subsection{Audio Tokenizers}

\textbf{DAC} (Descript Audio Codec)~\cite{kumar2024high} compresses 16\,kHz audio to 50\,Hz continuous latents with $d\!=\!1024$, followed by RVQ.
We use the pre-quantisation continuous latent as both source ($\mathbf{x}_0$) and target ($\mathbf{x}_1$).

\textbf{MOSS-Audio-Tokenizer}~\cite{gong2026moss} is a 1.6B-parameter purely Transformer-based audio tokenizer, trained end-to-end on 3M hours of diverse audio.
Its encoder (CAT architecture) consists of 4 hierarchical stages:
\begin{itemize}[itemsep=1pt,topsep=1pt]
  \item Stages 0--2: PatchedPretransform $\to$ 12-layer Transformer ($d\!=\!768$)
  \item Stage 3: PatchedPretransform $\to$ \textbf{32-layer Transformer} ($d\!=\!1280$)
\end{itemize}
Total temporal downsampling is $240 \!\times\! 2 \!\times\! 2 \!\times\! 2 = 1920$, yielding 12.5\,Hz features from 24\,kHz input.
We extract features from all 32 layers of Stage~3 for multi-layer conditioning.

% ============================================================
\section{Method}
\label{sec:method}

\subsection{System Overview}

Our pipeline consists of three stages (see \cref{fig:train,fig:inference}):

\textbf{(1) Offline feature extraction.}
All audio is processed once through DAC and MOSS before training:
DAC on clean speech $\to$ $\mathbf{x}_1 \in \mathbb{R}^{T \times 1024}$;
DAC on noisy speech $\to$ $\mathbf{x}_0 \in \mathbb{R}^{T \times 1024}$;
MOSS on noisy speech $\to$ conditioning features $\mathbf{c}$.
This offline design ensures that neither DAC (74M) nor MOSS (1.6B) are loaded during training, enabling training on Google Colab.

\textbf{(2) Training.}
At each step, we sample $t \sim U(0,1)$, compute $\mathbf{x}_t$ (\cref{eq:interpolation}), and train the DiT to predict the velocity field (\cref{eq:loss}).

\textbf{(3) Inference.}
Given a noisy utterance, we encode with DAC and MOSS, run a 50-step Euler ODE from $t\!=\!0$ to $t\!=\!1$, and decode the result with the DAC decoder.

\begin{figure}[t]
  \centering
  \includegraphics[width=\linewidth]{image1_train.png}
  \caption{\textbf{Training pipeline.} Clean and noisy speech are encoded offline by DAC and MOSS. The DiT predicts the velocity field $\mathbf{v}_\theta$ from the interpolated latent $\mathbf{x}_t$, conditioned on MOSS features $\mathbf{c}$ via cross-attention. Loss is MSE between predicted and true velocity.}
  \label{fig:train}
\end{figure}

\begin{figure}[t]
  \centering
  \includegraphics[width=\linewidth]{image2_inference.png}
  \caption{\textbf{Inference pipeline.} Starting from the noisy DAC latent $\mathbf{x}_0$, a 50-step Euler ODE solver iteratively refines the latent using the trained DiT. The result $\hat{\mathbf{x}}_1$ is decoded by the DAC decoder to produce enhanced speech.}
  \label{fig:inference}
\end{figure}

\subsection{Diffusion Transformer Architecture}

The velocity predictor is a DiT with 6 blocks ($\sim$28.4M parameters). Each block: (1)~AdaLN($t$) $\to$ Multi-Head Self-Attention (8 heads), (2)~AdaLN($t$) $\to$ Cross-Attention (Q from latent, K/V from $\mathbf{c}$), (3)~AdaLN($t$) $\to$ FFN with $4\times$ expansion.
Input/output projections map between DAC ($d\!=\!1024$) and hidden ($d\!=\!512$) dimensionality.

\subsection{Multi-Layer Conditioning}
\label{sec:conditioning}

We compare four conditioning strategies that differ \emph{only} in how MOSS features enter the cross-attention layers:

\textbf{Exp~1: No conditioning.}
Cross-attention layers removed. The DiT uses only the noisy latent and timestep. Baseline: 28.4M params.

\textbf{Exp~2: Last-layer conditioning.}
Standard approach: condition on $\mathbf{h}_{31} \in \mathbb{R}^{T_c \times 768}$, the final MOSS encoder output. Adds +10.1M params for cross-attention projections.

\textbf{Exp~3: Static multi-layer conditioning.}
All 32 hidden outputs $\{\mathbf{h}_i\}_{i=0}^{31} \in \mathbb{R}^{T_c \times 1280}$ are fused via:
\begin{equation}
  \mathbf{c} = \sum_{i=0}^{31} \text{softmax}(\mathbf{w})_i \cdot \mathbf{h}_i
  \label{eq:static}
\end{equation}
where $\mathbf{w} \in \mathbb{R}^{32}$ is learnable (initialised to uniform). Adds only 32 parameters over cross-attention.

\textbf{Exp~4: Time-dependent multi-layer conditioning.}
A small MLP ($1 \!\to\! 64 \!\xrightarrow{\text{SiLU}}\! 32$) maps the ODE timestep to per-layer weights:
\begin{equation}
  \mathbf{c}(t) = \sum_{i=0}^{31} \text{softmax}\big(\text{MLP}(t)\big)_i \cdot \mathbf{h}_i
  \label{eq:timedep}
\end{equation}
This adds 2,208 parameters and allows different denoising stages to emphasise different layers.

% ============================================================
\section{Experimental Setup}
\label{sec:setup}

\paragraph{Data.}
LibriSpeech dev-clean (2,703 utterances) mixed with MUSAN noise at SNR $\in \{-5, 0, 5, 10, 15\}$\,dB. All models trained on 5\,dB as proof-of-concept. Split: 2,163/270/270 (train/valid/test).

\paragraph{Features.}
DAC latents: $(T, 1024)$ truncated to 200 frames (4\,s at 50\,Hz). MOSS multi-layer: $32 \times (T_{\text{moss}}, 1280)$, truncated to 50 frames (12.5\,Hz).

\paragraph{Training.}
AdamW, lr $= 10^{-4}$, weight decay $10^{-5}$, batch 16, cosine schedule with 200-step warmup, max 100 epochs, early stopping (patience 20), gradient clipping at 1.0. Google Colab Pro (T4/A100).

\paragraph{Metrics.}
\textbf{PESQ} ($\uparrow$): perceptual speech quality~\cite{rix2001pesq}. \textbf{STOI} ($\uparrow$): short-time objective intelligibility~\cite{taal2010stoi}. \textbf{FAD} ($\downarrow$): Fr\'{e}chet Audio Distance in VGGish space~\cite{kilgour2019fad}.

% ============================================================
\section{Results \& Analysis}
\label{sec:results}

\subsection{Main Results}

\cref{tab:results} and \cref{fig:metrics} present the quantitative evaluation. All three metrics improve monotonically from Exp~1 through Exp~4.

\begin{table}[t]
\centering
\caption{Test set results (270 samples, 5\,dB SNR). Best in \textbf{bold}.}
\label{tab:results}
\small
\begin{tabular}{clccc}
\toprule
\textbf{Exp} & \textbf{Condition} & \textbf{PESQ $\uparrow$} & \textbf{STOI $\uparrow$} & \textbf{FAD $\downarrow$} \\
\midrule
1 & None            & 1.6048 & 0.8527 & 2.9774 \\
2 & Last layer      & 1.6499 & 0.8589 & 2.6997 \\
3 & Static ML       & 1.6868 & 0.8642 & 2.3857 \\
4 & Time-dep.\ ML   & \textbf{1.6986} & \textbf{0.8647} & \textbf{2.3456} \\
\bottomrule
\end{tabular}
\end{table}

\begin{figure}[t]
  \centering
  \includegraphics[width=\linewidth]{metrics_comparison.png}
  \caption{\textbf{Metrics comparison} across four conditioning variants. PESQ and STOI (higher is better) improve monotonically; FAD (lower is better) decreases consistently. Multi-layer variants (Exp~3, 4) substantially outperform single-layer (Exp~2) and unconditional (Exp~1) baselines.}
  \label{fig:metrics}
\end{figure}

\subsection{Analysis}

\paragraph{Multi-layer conditioning drives the largest gains.}
From Exp~1 to Exp~2: +2.8\% PESQ, $-$9.3\% FAD.
From Exp~2 to Exp~3: +5.1\% PESQ, $-$19.9\% FAD, with only 32 extra parameters.
This confirms that information across 32 MOSS layers is complementary.

\paragraph{FAD convergence (\cref{fig:fad}).}
This is perhaps our most revealing result. Exp~3 and Exp~4 both achieve $\sim$20\% lower FAD than Exp~1 at convergence. Notably, Exp~4 converges approximately 10 epochs faster than Exp~3, suggesting that time-dependent fusion provides a more efficient learning signal.
The faster convergence is practically significant: training on Colab with session time limits means that fewer epochs to convergence translates directly to more reliable experiment completion.

\begin{figure}[t]
  \centering
  \includegraphics[width=0.85\linewidth]{fad_over_epochs.png}
  \caption{\textbf{FAD convergence over training epochs.} Exp~3/4 (multi-layer) achieve $\sim$20\% lower FAD than Exp~1. Exp~4 (time-dependent) converges $\sim$10 epochs faster, indicating more efficient learning from adaptive layer weighting.}
  \label{fig:fad}
\end{figure}

\paragraph{Static fusion learns near-uniform weights (\cref{fig:fusion}).}
The learned static weight vector from Exp~3 has a Gini coefficient of only 0.007---nearly perfectly uniform.
All 32 layers contribute roughly equally, indicating that every layer adds non-redundant information.
This motivates the time-dependent variant: if all layers matter, dynamically re-weighting per timestep can only help.

\paragraph{Time-dependent weights reveal denoising dynamics (\cref{fig:timedep}).}
At $t \approx 0$ (near noisy input), the distribution is approximately uniform.
As $t \to 1$ (approaching clean target), the model progressively emphasises \emph{earlier layers} and the \emph{last layer}, with weight standard deviation increasing from 0.0012 to 0.0031.
This aligns with the intuition that early ODE steps benefit from broad contextual information, while later refinement steps require more specific acoustic detail (lower MOSS layers) and high-level coherence (last layer).

\begin{figure}[t]
  \centering
  \includegraphics[width=0.85\linewidth]{fusion_weights.png}
  \caption{\textbf{Static fusion weights} (Exp~3) across 32 MOSS layers. The distribution is nearly uniform (Gini\,=\,0.007), showing all layers contribute non-redundant conditioning information.}
  \label{fig:fusion}
\end{figure}

\begin{figure}[t]
  \centering
  \includegraphics[width=0.85\linewidth]{time_dependent_weights.png}
  \caption{\textbf{Time-dependent weights} (Exp~4) as a heatmap over ODE timesteps (y-axis) and MOSS layers (x-axis). Early timesteps show uniform weights; later timesteps concentrate on early layers and the last layer, adapting the conditioning to the denoising stage.}
  \label{fig:timedep}
\end{figure}

\paragraph{PESQ score context.}
PESQ of 1.60--1.70 may appear modest vs.\ supervised models.
However, flow-based generative models optimise naturalness over point-wise fidelity.
FAD, measuring distributional quality, shows 21.2\% improvement (Exp~1 $\to$ 4)---multi-layer conditioning most significantly improves perceptual naturalness.

% ============================================================
\section{Limitations \& Future Work}
\label{sec:limitations}

\textbf{Single-SNR training.}
All models were trained at 5\,dB only. Multi-SNR training is a natural extension.

\textbf{Limited noise diversity.}
MUSAN free-sound does not cover reverberation, babble, or music. Broader noise sources are needed for real-world robustness.

\textbf{32$\times$ feature storage.}
Storing 32 layers of $(T, 1280)$ is costly. A layer selection study (e.g., first 15 + last based on the weight analysis) could reduce this.

\textbf{DAC reconstruction ceiling.}
Output quality is bounded by DAC's encode-decode fidelity.

\textbf{Future directions} include: multi-SNR training, larger datasets (full LibriSpeech, DNS Challenge), ODE distillation~\cite{song2023consistency} for faster inference, and attention-based layer fusion.

% ============================================================
\section{Conclusion}

Through a controlled four-way ablation, we demonstrate that multi-layer conditioning from the MOSS-Audio-Tokenizer consistently outperforms both unconditional and last-layer-only baselines across PESQ, STOI, and FAD.
Our time-dependent fusion variant achieves the best results and converges faster, at a negligible parameter cost (+2,208).
The monotonic improvement---no conditioning $<$ last layer $<$ static multi-layer $<$ time-dependent---provides evidence that leveraging the full depth of pre-trained audio encoders benefits speech enhancement, analogous to multi-scale feature extraction in computer vision.

\paragraph{Acknowledgements.}
Training scripts for Google Colab were partially written with the assistance of Claude Opus (Anthropic).

% ============================================================
{\small
\bibliographystyle{ieee_fullname}
\begin{thebibliography}{10}

\bibitem{lipman2023flow}
Y.~Lipman, R.~T.~Q. Chen, H.~Ben-Hamu, M.~Nickel, and M.~Le.
\newblock Flow matching for generative modeling.
\newblock In \emph{ICLR}, 2023.

\bibitem{liu2023flow}
X.~Liu, C.~Gong, and Q.~Liu.
\newblock Flow straight and fast: Learning to generate and transfer data with rectified flow.
\newblock In \emph{ICLR}, 2023.

\bibitem{kumar2024high}
R.~Kumar, P.~Seetharaman, A.~Luebs, I.~Kumar, and K.~Kumar.
\newblock High-fidelity audio compression with improved {RVQGAN}.
\newblock In \emph{NeurIPS}, 2023.

\bibitem{gong2026moss}
Y.~Gong, K.~Chen, Z.~Fei, X.~Yang, K.~Chen, Y.~Wang, K.~Huang, M.~Chen, R.~Li, et~al.
\newblock {MOSS-Audio-Tokenizer}: Scaling audio tokenizers for future audio foundation models.
\newblock \emph{arXiv:2602.10934}, 2026.

\bibitem{peebles2023scalable}
W.~Peebles and S.~Xie.
\newblock Scalable diffusion models with transformers.
\newblock In \emph{ICCV}, 2023.

\bibitem{richter2023speech}
J.~Richter, S.~Welker, J.~Lemercier, B.~Lay, and T.~Gerkmann.
\newblock Speech enhancement and dereverberation with diffusion-based generative models.
\newblock \emph{IEEE/ACM TASLP}, 2023.

\bibitem{lu2013speech}
X.~Lu, Y.~Tsao, S.~Matsuda, and C.~Hori.
\newblock Speech enhancement based on deep denoising autoencoder.
\newblock In \emph{Interspeech}, 2013.

\bibitem{rix2001pesq}
A.~W. Rix, J.~G. Beerends, M.~P. Hollier, and A.~P. Hekstra.
\newblock Perceptual evaluation of speech quality ({PESQ}).
\newblock In \emph{ICASSP}, 2001.

\bibitem{taal2010stoi}
C.~H. Taal, R.~C. Hendriks, R.~Heusdens, and J.~Jensen.
\newblock A short-time objective intelligibility measure for time-frequency weighted noisy speech.
\newblock In \emph{ICASSP}, 2010.

\bibitem{kilgour2019fad}
K.~Kilgour, M.~Zuluaga, D.~Roblek, and M.~Sharifi.
\newblock Fr\'echet audio distance: A reference-free metric for evaluating music enhancement algorithms.
\newblock In \emph{Interspeech}, 2019.

\bibitem{song2023consistency}
Y.~Song, P.~Dhariwal, M.~Chen, and I.~Sutskever.
\newblock Consistency models.
\newblock In \emph{ICML}, 2023.

\end{thebibliography}
}

\end{document}
